\section{Diseño de sistemas multiagente para un co-científico de IA}

Los dos capítulos anteriores explicaron por qué se necesitan sistemas; este capítulo entra en el interior del sistema. Al leer el diagrama de arquitectura, debe seguirse tres flujos: cómo define el científico el problema, cómo los agentes hacen que las hipótesis candidatas mejoren continuamente y cómo las herramientas y la memoria conectan el ciclo interno de texto con la evidencia externa.

\subsection{La entrada del científico no es un prompt mágico universal}

El científico puede describir el objetivo de investigación, las restricciones experimentales, las preferencias atómicas y las ideas previas; también puede realizar add, review y discuss a las candidatas. El sistema devuelve finalmente hipótesis top-ranked, propuestas de investigación y una visión general. Esta interfaz reconoce que el propio objetivo cambia: tras ver los resultados intermedios, el experto puede corregir restricciones, señalar omisiones o redefinir el problema.

\lecturefigure{slide-007.jpg}{Arquitectura multiagente completa del co-científico de IA}{Grabación oficial de Stanford CS25 V6 Lecture 07 00:22:34; fuente arXiv:2502.18864v1}

\begin{knowledgebox}{Lectura de la figura: observe el control de izquierda a derecha, y el ciclo de mejora desde el centro}
A la izquierda, el científico define el objetivo y proporciona retroalimentación continua; en el centro, el supervisor configura el plan de investigación y el ranking tournament devuelve los resultados de comparación hacia generation, reflection, evolution y proximity; a la derecha, el sistema invoca herramientas como search y modelos especializados, y utiliza memory para mantener el estado entre rondas. La salida no es una respuesta aislada, sino un conjunto de objetos de investigación clasificados, con evidencia y sugerencias experimentales.
\end{knowledgebox}

\teachervoice{00:22:26--00:25:54, el docente enfatiza que el científico puede proporcionar seed idea, restricciones y retroalimentación en lenguaje natural. La "autonomía" del sistema es una autonomía de exploración controlada por objetivos y recursos, no un poder de investigación delegado independiente del experto.}

\begin{importantbox}{Scientist-in-the-loop no equivale a firmar al final}
El verdadero scientist-in-the-loop debe cubrir la definición del problema, la selección de evidencia, el feedback a las candidatas, el diseño experimental, la interpretación de resultados y las condiciones de parada. Si la persona solo hace clic en "Aceptar" después de que el sistema genera un informe de cien páginas, su rol se ha degradado de colaborador al mero responsable, incapaz de corregir la dirección de la búsqueda.
\end{importantbox}

\subsection{Divisiones de trabajo entre las seis categorías centrales de Agentes}

Los agentes en la arquitectura no son roles personificados, sino encapsulamientos modulares de diferentes operaciones de razonamiento. Generation se encarga de expandir el espacio de candidatas; Reflection utiliza literatura, simulaciones y contraejemplos para reducir errores evidentes; Ranking realiza comparaciones relativas; Evolution combina y reescribe candidatas de alto valor; Proximity juzga las relaciones entre candidatas; y Meta-review extrae problemas comunes y pasos siguientes del conjunto completo.

\begin{knowledgebox}{Glosario: por qué los Agentes deben dividirse el trabajo}
\begin{tabular}{L{0.17\textwidth}L{0.33\textwidth}L{0.32\textwidth}}
\toprule
Agente & Operación central & Riesgo principal de falla \\
\midrule
Generation & Exploración bibliográfica, debate científico, formulación de hipótesis diversas & Repetición de conclusiones conocidas, fantasías sin restricciones \\
Reflection & Review completo, web search, simulation review, deep verification & Omisión de evidencia, rechazo erróneo de ideas novedosas \\
Ranking & Comparaciones por pares, torneos, resumen de patrones win/loss & Preferencias del judge y evaluaciones circulares \\
Evolution & Combinación, simplificación y expansión de candidatas de alto valor & Solo pulido estilístico sin añadir contenido verificable \\
Proximity & Juicio de similitud, complementariedad o conflicto entre candidatas & Similitud superficial semántica que oculta diferencias mecánicas \\
Meta-review & Formación de panorama de investigación y conclusiones transversales & Pérdida de pistas clave pero minoritarias al comprimir \\
\bottomrule
\end{tabular}
\end{knowledgebox}

\begin{lstlisting}[language=Python,caption={Ciclo conceptual Generate--debate--evolve}]
def research_cycle(goal, constraints, budget):
    hypotheses = generation.explore(goal, constraints)
    while budget.remaining():
        reviews = reflection.critique(hypotheses, tools=True)
        matches = ranking.pairwise_tournament(hypotheses, reviews)
        ranked = ranking.update_scores(matches)
        hypotheses = evolution.expand(ranked.top(), reviews)
        hypotheses = proximity.merge_and_separate(hypotheses)
        memory.store(hypotheses, reviews, matches)
    return meta_review.synthesize(hypotheses, memory)
\end{lstlisting}

\subsection{Ranking tournament y Elo: por qué las comparaciones relativas son útiles}

Las hipótesis abiertas son difíciles de puntuar directamente con un valor absoluto, pero entre dos candidatas se pueden plantear preguntas más específicas: ¿cuál tiene un mecanismo más completo? ¿Cuál es más falsable? ¿Cuál entra en menos conflicto con la evidencia existente? Por ello, el agente Ranking organiza un pairwise tournament y actualiza las prioridades según los patrones win/loss.

Si las puntuaciones actuales de las candidatas $i$ y $j$ son $r_i,r_j$, la probabilidad de victoria esperada estilo Elo se puede escribir como
$$
E_i=\frac{1}{1+10^{(r_j-r_i)/400}},
$$
actualizando a
$$
r_i' = r_i + \eta(S_i-E_i).
$$
Donde:
\begin{itemize}
  \item $E_i$ es la probabilidad de victoria de la candidata $i$ predicha por el sistema según las puntuaciones anteriores;
  \item $S_i\in\{0,0.5,1\}$ es el resultado de esta comparación;
  \item $\eta$ es el tamaño del paso de actualización;
  \item $r_i'$ solo representa la prioridad relativa dentro del judge actual, rubric y conjunto de candidatas.
\end{itemize}

\begin{warningbox}{Elo no puede probar novelty, truth o impact}
Si el judge prefiere hipótesis con terminología densa, estructura completa o cercanas a la distribución de entrenamiento, Elo amplificará estabilmente esta preferencia. El ordenamiento científico debe registrar dimensiones de evaluación, razones de comparación, evidencia citada y contraejemplos; finalmente, se requiere que expertos y experimentos conecten el ranking interno con el mundo real.
\end{warningbox}

\teachervoice{00:29:12--00:31:47, el docente compara ranking con torneos de Go o ajedrez: lo importante es hacer que las candidatas se comparen entre sí y resuman patrones de fracaso. Esta analogía explica los mecanismos de ingeniería, pero no implica que la ciencia posea victorias/derrotas tan objetivas e inmediatas como en los juegos de tablero.}

\subsection{Test-time compute y Supervisor asíncrono}

Test-time compute se refiere al cómputo adicional de inferencia, muestreo, búsqueda, herramientas y verificación invertido durante una tarea de investigación una vez que los parámetros del modelo están fijos. El supervisor del co-científico de IA desglosa objetivos de alto nivel en tareas, las escribe en colas asíncronas y asigna diferentes modelos y presupuestos según el valor de la candidata y las necesidades de verificación. La asincronía permite que generation no espere a que todos los reflection se completen, y también permite que verificaciones costosas sirvan solo a un pequeño número de candidatas de alta prioridad.

\begin{knowledgebox}{Concepto previo: test-time compute}
El cómputo de entrenamiento cambia los parámetros; el test-time compute no actualiza los pesos del modelo base, sino que aumenta la profundidad de búsqueda, el número de candidatas, las rondas de reflexión y las llamadas a herramientas para el problema actual. Puede mejorar la calidad en tareas específicas, pero simultáneamente incrementa el consumo de tokens, latencia, costos de APIs externas y de revisión; debe gestionarse junto con el presupuesto experimental.
\end{knowledgebox}

El presupuesto total $B$ puede ser consumido por diferentes tareas:
$$
\sum_{a\in\mathcal{A}} c_a n_a
+\sum_{t\in\mathcal{T}} c_t n_t
+C_{\text{human}}
\le B.
$$
Donde:
\begin{itemize}
  \item $\mathcal{A}$ es el conjunto de tipos de agente, $n_a$ es el número de llamadas y $c_a$ es el costo unitario;
  \item $\mathcal{T}$ es el conjunto de herramientas externas, $n_t$ y $c_t$ representan respectivamente el número de llamadas y el costo;
  \item $C_{\text{human}}$ es el costo de lectura, feedback y diseño experimental del experto;
  \item $B$ no puede ser solo dólares, sino también latencia, tiempo de GPU o capacidad experimental.
\end{itemize}

\begin{lstlisting}[language=Python,caption={Esquema de enrutamiento de recursos del Supervisor asíncrono}]
while queue.has_work() and budget.available():
    task = queue.pop_highest_value()
    model = "pro" if task.requires_deep_verification else "flash"
    result = workers.submit(task, model=model)
    memory.link(result, task.provenance)
    queue.add(reflection.followups(result))
    queue.add(evolution.promising_variants(result))
\end{lstlisting}

\teachervoice{00:39:11--00:39:43, el docente ofrece experiencia de enrutamiento real: review o verification pueden usar la variante Pro con mayor capacidad de razonamiento, mientras que trabajos más simples usan Flash para ahorrar tokens y costos. Las diapositivas lo tratan como workload-aware routing, no como una recomendación de producto fija.}

\subsection{Ciclicidad de la evaluación interna}

Una pregunta planteada a mitad del curso golpeó el núcleo del problema: si modelos del mismo tipo generan hipótesis y luego modelos del mismo tipo las evalúan, ¿el sistema solo recompensará textos que "suenen como buena ciencia"? Los papers públicos usan objetivos de expertos, Elo automático y comparaciones multi-modelo para demostrar test-time scaling, pero estas evidencias muestran principalmente que el flujo interno puede optimizar continuamente cierta rubric, sin equivaler a que el nuevo mecanismo se establezca en la naturaleza.

\begin{warningbox}{El triple ciclo del Model-as-judge}
El generador y el judge pueden compartir preferencias de datos de entrenamiento; un novelty judge podría no conocer investigaciones no publicadas; una narrativa completa puede aplastar pistas verdaderamente clave pero mal expresadas. Reducir la ciclicidad requiere jueces heterogéneos, rubrics claras, recuperación de evidencia, evaluación ciega por expertos, segmentación temporal y experimentos finales, en lugar de añadir otra ronda de "por favor revisa con cuidado".
\end{warningbox}

\teachervoice{00:33:58--00:37:04, el docente admite que el sistema no gana en todos los objetivos; también la audiencia cuestionó si la evaluación automática puede medir realmente la novelty. El curso no ocultó estas preguntas detrás de promedios.}

\subsection{Resumen de la sección}

El co-científico de IA organiza la entrada del científico, seis tipos de Agentes, Supervisor, herramientas y memoria en un ciclo generate--debate--rank--evolve. Los pairwise tournament son adecuados para proporcionar prioridades relativas a candidatas abiertas; el test-time compute permite aumentar la búsqueda según la tarea; pero Elo y model-as-judge siguen siendo solo señales internas. La fiabilidad del sistema depende de poder conectar la búsqueda interna con evidencia heterogénea, juicio humano y experimentos reales.


